\documentclass[12pt]{article}
\usepackage[utf8]{inputenc}
\usepackage[T1]{fontenc}
\usepackage[margin=1in]{geometry}
\usepackage{setspace}
\setlength{\parindent}{2em}
\setlength{\parskip}{1.25em}
\renewcommand{\baselinestretch}{1.0}
\makeatletter
\newcommand{\@subtitle}{}
\newcommand{\subtitle}[1]{\renewcommand{\@subtitle}{#1}}
\renewcommand{\@maketitle}{%
  \newpage
  \null
  \vskip 1em%
  \begin{center}%
    {\LARGE \@title \par}%
    \ifx\@subtitle\@empty\else
      \vskip 0.5em%
      {\large \@subtitle \par}%
    \fi
  \end{center}%
  \par
  \vskip 1em}
\makeatother
\usepackage{xcolor}
\usepackage[colorlinks=true,urlcolor=blue!55!black,linkcolor=black]{hyperref}

\title{\textit{City of Memphis v. Greene}, 451 U.S. 100 (1981)}
\subtitle{Case Brief}
\author{}
\date{}

\begin{document}
\maketitle

\section{Memory Jogger}

A city's closure of a residential street between an all-white neighborhood and a predominantly Black neighborhood to reduce through traffic did not violate 42 U.S.C.\ \S~1982 or the Thirteenth Amendment, because it caused only a routine commuting inconvenience without discriminatory intent or impairment of Black property rights (\href{https://www.courtlistener.com/opinion/110461/city-of-memphis-v-greene/}{CourtListener}; \href{https://supreme.justia.com/cases/federal/us/451/100/}{Justia}).

\section{Detailed Case Facts}

Hein Park was a small residential subdivision in Memphis, Tennessee, developed before World War II as an exclusive white neighborhood; at the time of the litigation, every home in Hein Park was owned by white residents. The area immediately north of Hein Park was predominantly Black. West Drive, a two-lane street roughly half a mile long, ran north--south through the center of Hein Park. Its northern terminus met the intersection of Jackson Avenue and Springdale Street (two heavily traveled four-lane avenues), while its southern terminus ended near Overton Park, a large 342-acre municipal park containing the Memphis Zoo and other public amenities. West Drive was one of three streets entering Hein Park from the north, with two additional streets entering from the east (\href{https://supreme.justia.com/cases/federal/us/451/100/}{451 U.S. at 102--103 \& n.2}).

In 1970, Hein Park residents petitioned Memphis to close four streets leading into the subdivision, but the city denied that request after police, fire, and sanitation objections. However, the Traffic Engineering Department noted that closing West Drive at Jackson Avenue would eliminate much of the through traffic. On July 9, 1973, members of the Hein Park Civic Association filed an ``Application to Close Streets or Alleys'' with the Memphis and Shelby County Planning Commission to close the northernmost 25 feet of West Drive just south of Jackson Avenue and deed that strip to the abutting property owners (leaving a 24-foot gap with a speed breaker for municipal and emergency vehicles). The application was signed by the two property owners abutting both Jackson Avenue and West Drive and all but one homeowner (Sarah Terry) on that block (\href{https://supreme.justia.com/cases/federal/us/451/100/}{451 U.S. at 103--104, 112 n.19, 115 n.26}).

The applicants stated three reasons for the closure: (1) reducing through traffic on subdivision streets; (2) increasing safety for children living in Hein Park and walking to Snowden Junior High School; and (3) reducing ``traffic pollution'' such as noise, litter, and interruption of community living. On November 1, 1973, the Planning Commission recommended approval subject to utility easements and fire-vehicle clearance. Despite opposition from residents north of Jackson Avenue---including petitions bearing roughly 1,000 signatures protesting that the closure symbolized a white neighborhood shutting its door on adjacent Black communities---the Memphis City Council held hearings and approved the closure resolution, later reconsidering but never rescinding it (\href{https://supreme.justia.com/cases/federal/us/451/100/}{451 U.S. at 104--105 \& n.5, 143 n.9}).

The dispute arose in a city with a documented history of official resistance to desegregation (see \textit{Watson v. Memphis}, 373 U.S. 526 (1963)) and tested the post-Reconstruction reach of two federal provisions: Section~1 of the Civil Rights Act of 1866, codified at 42 U.S.C.\ \S~1982 (guaranteeing all citizens ``the same right, in every State and Territory, as is enjoyed by white citizens thereof to inherit, purchase, lease, sell, hold, and convey real and personal property''), and the Thirteenth Amendment, which abolishes slavery and involuntary servitude and empowers Congress in Section~2 to enforce that ban by appropriate legislation (\href{https://supreme.justia.com/cases/federal/us/451/100/}{451 U.S. at 102 \& n.1, 116 n.27}).

\section{Procedural History}

On April 1, 1974, three individuals (including N. T. Greene) and two civic associations filed a class action in the United States District Court for the Western District of Tennessee against the City of Memphis and city officials on behalf of Black property owners north of Jackson Avenue and west of Springdale Street, seeking an injunction under 42 U.S.C.\ \S\S~1982 and 1983 and the Thirteenth and Fourteenth Amendments. The District Court initially dismissed the complaint for failure to state a claim, finding no injury to plaintiffs' own property, no disparate racial effect, and no standing to raise procedural objections (\href{https://supreme.justia.com/cases/federal/us/451/100/}{451 U.S. at 105--106}).

The United States Court of Appeals for the Sixth Circuit reversed and remanded, 535 F.2d 976 (6th Cir.\ 1976), holding that plaintiffs stated a viable claim under \S~1982 or \S~1983 if they could prove that city officials granted the privacy and quiet of a dead-end street to white residents ``because of their color'' while denying comparable benefits to Black residents. On remand, District Judge Robert M. McRae certified a Rule 23(b)(2) class, conducted a full trial, and entered judgment for the city. Although Judge McRae found that the closure would have a disproportionate impact on ``overwhelmingly black'' southbound traffic, he concluded that the city had not denied Black citizens a benefit granted to white citizens, had not acted with racially discriminatory intent or purpose under \textit{Arlington Heights v. Metropolitan Housing Development Corp.}, and had not departed significantly from normal procedures (\href{https://supreme.justia.com/cases/federal/us/451/100/}{451 U.S. at 106--108 \& nn.10--12}).

On the second appeal, the Sixth Circuit again reversed and remanded, 610 F.2d 395 (6th Cir.\ 1979). Without overturning Judge McRae's factual findings, the appellate panel held that plaintiffs did not need to prove unequal denial of street-closing applications; rather, the unique closure at the racial dividing line---combined with testimony about ``undesirable'' traffic and potential property-value depreciation---constituted a ``badge of slavery'' impairing Black property rights in violation of \S~1982 and the Thirteenth Amendment. The Supreme Court granted certiorari (446 U.S. 934 (1980)), heard oral argument on December 3, 1980, and reversed on April 20, 1981 (\href{https://supreme.justia.com/cases/federal/us/451/100/}{451 U.S. at 108--110, 129}).

\section{Judicial Votes}

\begin{itemize}
  \item \textbf{Majority (6--3 judgment; 5-Justice opinion of the Court):} Justice John Paul Stevens delivered the opinion of the Court, joined by Chief Justice Warren E. Burger and Justices Potter Stewart, Lewis F. Powell, Jr., and William H. Rehnquist.
  \item \textbf{Concurrence in the Judgment:} Justice Byron R. White filed an opinion concurring in the judgment.
  \item \textbf{Dissent:} Justice Thurgood Marshall filed a dissenting opinion, joined by Justices William J. Brennan, Jr., and Harry A. Blackmun.
\end{itemize}

\section{Holding}

\textbf{Judgment:} Reversed. The Sixth Circuit's judgment invalidating the closure of West Drive is reversed (\href{https://supreme.justia.com/cases/federal/us/451/100/}{451 U.S. at 129}).

\textbf{Rule Announced:} A facially neutral municipal street closure motivated by legitimate traffic-safety and residential-tranquility interests---without proof of racially discriminatory intent, property-value depreciation, severe restriction of access to homes, or unequal denial of municipal benefits---does not impair property rights protected by 42 U.S.C.\ \S~1982 and does not constitute a ``badge or incident of slavery'' violating Section~1 of the Thirteenth Amendment, even if it disparately inconveniences Black motorists (\href{https://supreme.justia.com/cases/federal/us/451/100/}{451 U.S. at 123--129}).

\textbf{Holding vs.\ Dicta:} The Court's holding is narrow and record-bound: on the District Court's factual findings, the minor rerouting of traffic neither implicated \S~1982 property interests nor rose to a Thirteenth Amendment badge of slavery. Crucially, the majority explicitly left open two major doctrinal questions as undecided (not holding): (1) whether a violation of 42 U.S.C.\ \S~1982 requires proof of purposeful racial discrimination, and (2) whether Section~1 of the Thirteenth Amendment, by its own unaided force without Section~2 enabling legislation, prohibits ``badges and incidents of slavery'' beyond literal chattel slavery or involuntary servitude (\href{https://supreme.justia.com/cases/federal/us/451/100/}{451 U.S. at 120, 125--126, 128--129}).

\section{Analysis of Opinions}

\subsection{Majority Opinion (Justice Stevens)}

Justice Stevens first audited the four factual predicates on which the Sixth Circuit relied and concluded that the trial record did not support them. While accepting Judge McRae's finding that diverted southbound motorists were predominantly Black, Stevens characterized the actual inconvenience as minimal because West Drive was only a half-mile two-lane street dead-ending near Overton Park; southbound drivers merely had to turn east or west a few blocks earlier. Second, the barrier was only a curb with a 24-foot gap for municipal vehicles, and other streets still provided access to Hein Park. Third, the isolated hearsay remark by an unnamed petition solicitor calling the traffic ``undesirable traffic'' did not establish racial animus by the City Council, which acted on legitimate concerns over traffic volume (1,600 to 1,700 cars per 12-hour period) and child safety near Snowden Junior High School. Fourth, plaintiffs' own real estate broker testified that property values north of Hein Park would not decline economically, offering only uncredited speculation about future psychological demoralization (\href{https://supreme.justia.com/cases/federal/us/451/100/}{451 U.S. at 110--119}).

Turning to the law, Justice Stevens noted that under \textit{Washington v. Davis}, 426 U.S. 229 (1976), and \textit{Arlington Heights}, 429 U.S. 252 (1977), the absence of discriminatory intent foreclosed any Fourteenth Amendment Equal Protection claim. Rather than deciding whether \S~1982 or the Thirteenth Amendment similarly requires proof of discriminatory intent, Stevens resolved both claims on threshold scope grounds. Reviewing \textit{Hurd v. Hodge}, \textit{Jones v. Alfred H. Mayer Co.}, \textit{Sullivan v. Little Hunting Park, Inc.}, and \textit{Tillman v. Wheaton-Haven Recreation Ass'n}, he held that \S~1982 protects the equal right to acquire, hold, and use property without racial impairment. Because Memphis did not deny Black neighborhoods comparable closures, did not depreciate Black property values, and did not severely restrict access to Black homes (distinguishing \textit{Jennings v. Patterson}, 488 F.2d 436 (5th Cir.\ 1974)), requiring motorists to use an alternate public street did not impair a \S~1982 property interest (\href{https://supreme.justia.com/cases/federal/us/451/100/}{451 U.S. at 119--124}).

On the Thirteenth Amendment claim, Justice Stevens reserved the question left open in \textit{Jones}---whether Section~1 independently bans ``badges and incidents of slavery''---because the closure could not fairly be characterized as one. Local governments enjoy wide discretion to balance motorist convenience against pedestrian safety and residential tranquility (\textit{Village of Belle Terre v. Boraas}). Because urban neighborhoods are frequently racially concentrated, ordinary traffic regulations inevitably produce disparate impacts; treating such ``routine burdens of citizenship'' as a Thirteenth Amendment stigma would ``trivialize the great purpose of that charter of freedom'' (\href{https://supreme.justia.com/cases/federal/us/451/100/}{451 U.S. at 124--129}).

\subsection{Opinion Concurring in the Judgment (Justice White)}

Justice White refused to join the majority's reexamination of the trial record, rebuking the Court for disregarding Supreme Court Rule~21.1(a) by acting as a factfinder on unbriefed issues rather than answering the sole question presented in the certiorari petition: whether a nonintentional disparate impact violates 42 U.S.C.\ \S~1982. Answering that question directly, White examined the text and legislative history of the Civil Rights Act of 1866, enacted under Section~2 of the Thirteenth Amendment to dismantle the post-Civil War Black Codes and direct, intentional oppression of freedmen (citing Senator Trumbull's floor remarks and Major General Carl Schurz's 1865 report). Concluding that nothing in the 1866 debates suggested Congress targeted facially neutral laws with incidental racial effects, White held that \S~1982 requires proof of racial animus or purposeful racial discrimination. Because the Sixth Circuit reversed without disturbing the District Court's finding of no discriminatory intent, White concurred in the judgment of reversal (\href{https://supreme.justia.com/cases/federal/us/451/100/}{451 U.S. at 129--135}).

\subsection{Dissenting Opinion (Justice Marshall)}

Justice Marshall, joined by Justices Brennan and Blackmun, accused the majority of sanitizing a textbook racial barrier into a benign traffic measure. Emphasizing that Hein Park was developed as an all-white enclave and that the barrier sat precisely at the boundary separating it from a predominantly Black community, Marshall argued that the closure sent an unmistakable stigmatizing message---what expert psychiatrist Dr. Marvin Feit called a ``monument to racial hostility''---telling Black citizens to stay out of a protected white neighborhood. Moreover, deeding the 25-foot strip to private abutting owners empowered them to exclude pedestrians as well as vehicles (\href{https://supreme.justia.com/cases/federal/us/451/100/}{451 U.S. at 135--140}).

Marshall marshaled substantial \textit{Arlington Heights} circumstantial evidence of discriminatory purpose ignored by the majority: Sarah Terry's testimony regarding ``undesirable traffic''; the city's procedural deviations (processing the application without all abutting owners' signatures, giving no notice to Black residents north of Jackson Avenue, denying access to the Planning Commission file, and limiting opponents to 15 minutes before the City Council); the fact that Memphis had never before closed a through street for traffic control; and Memphis's documented history of segregation. He also credited real estate broker H. C. Moore's and Dr. Feit's uncontroverted testimony predicting long-term property depreciation and heightened police harassment north of Hein Park (\href{https://supreme.justia.com/cases/federal/us/451/100/}{451 U.S. at 141--147}). Construing \S~1982 broadly per \textit{Jones} and the 1866 debates (which expressly repudiated Senator Van Winkle's defense of white racial enclaves), Marshall argued that official action causing known, massively disproportionate racial harm to property enjoyment required a heavy government justification that Memphis failed to meet. Finally, in Footnote~18, Marshall argued that because \S~1982 merely codified the guarantees implicit in Section~1 of the Thirteenth Amendment, the barrier violated the Thirteenth Amendment as a badge or incident of slavery as well (\href{https://supreme.justia.com/cases/federal/us/451/100/}{451 U.S. at 147--155 \& n.18}).

\section{Examples: Future Applications}

\begin{enumerate}
  \item \textbf{Same-Side (Upheld under \textit{Greene}):} A city converts a two-lane residential street running through a predominantly white historic district into a one-way street with speed humps to curb rush-hour cut-through traffic from an adjacent predominantly Hispanic neighborhood, where commuters retain direct access to parallel four-lane arterial roads. Because the measure serves legitimate traffic-calming goals, is not motivated by racial animus, and does not depress neighboring property values or block home access, it is a ``routine burden of citizenship'' valid under \S~1982 and the Thirteenth Amendment.
  \item \textbf{Same-Side (Upheld under \textit{Greene}):} A county relocates a regional bus transfer hub half a mile south for transit-efficiency reasons, adding four minutes to the commute of residents of a predominantly Black subdivision without proof of discriminatory purpose or impairment of real-property ownership or value. Under \textit{Greene}, indirect commuting inconvenience lacking a nexus to property rights fails to state a claim under \S~1982 or Section~1 of the Thirteenth Amendment.
  \item \textbf{Opposite-Side (Invalidated under \textit{Greene}'s Distinguishing Criteria):} A town erects a permanent locked gate across the only paved road connecting a predominantly Black residential enclave to the municipal center, forcing Black homeowners to drive an extra two miles over an unpaved detour while issuing gate keys to the lone white landowner on the enclave's side. As the \textit{Greene} majority noted in Footnote~36 (distinguishing \textit{Jennings v. Patterson}), severely restricting physical ingress and egress to Black-owned homes while granting white owners access directly impairs property rights under \S~1982.
  \item \textbf{Opposite-Side (Invalidated under \textit{Greene}'s Distinguishing Criteria):} A city approves cul-de-sac street closures for five white subdivisions upon petition, but repeatedly denies identical traffic-safety closure petitions submitted by three similarly situated Black neighborhoods despite comparable traffic volumes and accident records. As \textit{Greene} explicitly acknowledged (\href{https://supreme.justia.com/cases/federal/us/451/100/}{451 U.S. at 123}), conferring property-enhancing municipal benefits on white homeowners while refusing the same benefits to Black homeowners violates \S~1982.
  \item \textbf{Fence-Sitter (Close Case):} A city builds a concrete sound wall with no pedestrian openings along the boundary between an affluent white subdivision and a newly built low-income, predominantly Black apartment complex, citing highway noise abatement while also responding to residents' complaints about ``foot traffic and crime'' from the apartments. Unlike \textit{Greene}, the wall completely severs pedestrian passage and carries stronger evidence of pretextual, race-coded motive; yet the city will argue under \textit{Greene} that acoustic buffering is a legitimate municipal objective and that apartment tenants still retain vehicular street access on the other side of the complex.
\end{enumerate}

\section{Critique}

\textit{City of Memphis v. Greene} has been widely criticized for hollowing out the Thirteenth Amendment and narrowing 42 U.S.C.\ \S~1982 just thirteen years after the Court's expansive reading in \textit{Jones v. Alfred H. Mayer Co.} (1968). Structurally, the majority opinion suffered from a sharp methodological flaw identified by Justice White in concurrence: by ducking the certified legal question (whether \S~1982 requires discriminatory intent) and instead reweighing trial evidence to minimize the barrier's practical and psychological harm, the Court left both \S~1982 intent doctrine and Section~1 Thirteenth Amendment doctrine in prolonged limbo (\href{https://dsc.duq.edu/cgi/viewcontent.cgi?article=2263&context=dlr}{Duquesne Law Review}).

From a \textbf{progressive and antisubordination perspective}, scholars argue that Justice Stevens adopted a formalistic, ahistorical lens that echoed \textit{Plessy v. Ferguson}'s dismissal of racial stigma as merely ``symbolic'' (\href{https://arizonalawreview.org/pdf/62-3/62arizlrev637.pdf}{Michael A. Lawrence, 62 Ariz.\ L. Rev.\ 637, 660 (2020)}). By insisting that Black motorists' inconvenience was merely a function of ``where they live'' rather than race, the majority erased the legally enforced history of residential segregation in Memphis that created all-white Hein Park in the first place. Furthermore, by reserving whether Section~1 of the Thirteenth Amendment independently prohibits badges and incidents of slavery without providing any judicial standard, \textit{Greene} effectively froze Section~1 litigation in the lower federal courts, which subsequently treated the Amendment as reaching only literal coerced labor (\href{https://lawreview.law.ucdavis.edu/sites/g/files/dgvnsk15026/files/media/documents/DavisVol40No4_Carter.pdf}{William M. Carter, Jr., 40 U.C. Davis L. Rev.\ 1311 (2007)}).

From an \textbf{originalist and textualist perspective}, the critique cuts in two directions. On one hand, textualists generally applaud the Court's refusal to transform a localized traffic barrier into ``slavery'' or ``involuntary servitude'' under Section~1 of the Thirteenth Amendment, as well as Justice White's historically grounded reading of the 1866 Civil Rights Act as targeting purposeful racial deprivations rather than incidental disparate impacts (\href{https://constitution.congress.gov/browse/essay/amdt13-S2-3/ALDE_00013216/}{Constitution Annotated}). From this standpoint, stretching Section~1 to allow federal courts---without congressional authorization under Section~2---to invalidate facially neutral local traffic rules would violate separation of powers and federalism. On the other hand, some Reconstruction-era originalists contend that the 38th and 39th Congresses understood Section~1 itself to obliterate core civil disabilities and racially exclusionary customs, not merely chattel bondage, and fault \textit{Greene} for leaving Section~1's original public meaning unresolved while deferring uncritically to municipal pretexts (\href{https://www.uclalawreview.org/wp-content/uploads/2019/09/Pope-65.2.pdf}{James Gray Pope, 65 UCLA L. Rev.\ 1466 (2018)}).

\section{Key Quotations}

\begin{quote}
The injury to respondents established by the record is the requirement that one public street rather than another must be used for certain trips within the city. We need not assess the magnitude of that injury to conclude that it does not involve any impairment to the kind of property interests that we have identified as being within the reach of \S~1982.
--- Justice Stevens, Majority Opinion, \href{https://supreme.justia.com/cases/federal/us/451/100/}{451 U.S. at 124}.
\end{quote}

\begin{quote}
Because urban neighborhoods are so frequently characterized by a common ethnic or racial heritage, a regulation's adverse impact on a particular neighborhood will often have a disparate effect on an identifiable ethnic or racial group. To regard an inevitable consequence of that kind as a form of stigma so severe as to violate the Thirteenth Amendment would trivialize the great purpose of that charter of freedom.
--- Justice Stevens, Majority Opinion, \href{https://supreme.justia.com/cases/federal/us/451/100/}{451 U.S. at 128}.
\end{quote}

\begin{quote}
We merely hold that the impact of the closing of West Drive on nonresidents of Hein Park is a routine burden of citizenship; it does not reflect a violation of the Thirteenth Amendment.
--- Justice Stevens, Majority Opinion, \href{https://supreme.justia.com/cases/federal/us/451/100/}{451 U.S. at 129}.
\end{quote}

\begin{quote}
We have no basis for concluding anything other than that a violation of \S~1982 requires some showing of racial animus or an intent to discriminate on the basis of race.
--- Justice White, Concurring in the Judgment, \href{https://supreme.justia.com/cases/federal/us/451/100/}{451 U.S. at 135}.
\end{quote}

\begin{quote}
Petitioner city of Memphis, acting at the behest of white property owners, has closed the main thoroughfare between an all-white enclave and a predominantly Negro area of the city. The stated explanation for the closing is of a sort all too familiar: ``protecting the safety and tranquility of a residential neighborhood'' by preventing ``undesirable traffic'' from entering it. Too often in our Nation's history, statements such as these have been little more than code phrases for racial discrimination. These words may still signify racial discrimination, but apparently not, after today's decision, forbidden discrimination.
--- Justice Marshall, Dissenting Opinion, \href{https://supreme.justia.com/cases/federal/us/451/100/}{451 U.S. at 135--136}.
\end{quote}

\end{document}
